\cotefile{157-4}

\scanpage{212}

\mspage{86}

\uline{Théorème 7.1} Soit $M$ une catégorie stable par petites \uline{$\lim$} filtrantes, \deleted{\unsure{engendrée|une partie}} $M_0\subset\operatorname{Ob}M$ un \inserted{\deleted{\unsure{e}}sous-}ensemble de $\operatorname{Ob}M$ générateurs par monomorphismes stricts, plus précisément on suppose \deleted{que pour} ceci.

(7.1.1) Pour \deleted{tout} monomorphisme $i\colon A\hookrightarrow B$ dans $M$ qui n'est pas iso, il existe un monomorphisme $B'\overset{j}{\hookrightarrow}B$ \unsure{ayant} les propriétés suivantes :

a) $B'\in M_0$

b) $j$ ne se factorise pas par $i$

c) $A'=A\times_B B'$ existe dans $M$\inserted{ et est dans $M_0$}, ainsi que \deleted{\unsure{$C$}} $C=A\amalg_{A'}B'$, et \deleted{la flèche} dans le diagramme canonique

\diagram{\begin{tikzcd}[ampersand replacement=\&,column sep=large]
 A' \arrow[r,hook,"i'"] \arrow[d,hook] \& B' \arrow[d,hook] \arrow[dr,hook,dashed,"j"] \& {}\\
 A \arrow[r,hook,"\alpha"] \arrow[rr,hook,dashed,bend right=19,"i"'] \& C \arrow[r,"\beta"] \& B
\end{tikzcd}}

la flèche $\beta$ est mono.

[\uline{NB} Ces conditions sont satisfaites si on suppose, en plus de la condition que $M_0$ soit \unsure{gén.} pour les monos, que $M_0$ contient, avec tt objet\inserted{ $B'$}, et \deleted{tout} \inserted{pour} \deleted{\unsure{partie}} objet $A'$ qui est \deleted{\unsure{isom.}} soit \deleted{\unsure{et}} un quotient, \deleted{\unsure{soit}} un sous-objet de $B'$, un objet $A'_0$ isom. à $A'$; et si de plus $M$ satisfait aux conditions d'exactitude suivantes : $1^\circ)$ \deleted{\unsure{toute}} flèche $f$ de $M$ se factorise en $ip$, avec $i$ mono, $p$ épi, et $2^\circ)$ Si $A_1,A_2$ sont deux sous-objets d'un objet $B$ de $M$, alors l'intersection $A_0=A_1\times_A A_2$ existe, ainsi que \unsure{la|le}
% archive sheet number: 211
\scanpage{213}

\mspage{87}

somme amalgamée $C=A_1\amalg_{A_0} A_2$, et le morphisme canonique $C\longrightarrow B$ est \deleted{un} mono.] \inserted{Supposons de plus que $M$ satisfait à la condition \textcircled{e} p.~76, i.e. \unsure{une|que} \deleted{les} $\varinjlim$ filtr. de monos est mono.} \deleted{Alors \unsure{on a ce qui suit}}

$1^\circ)$ Soit $\mathrm{Mon}\subset\mathrm{Fl}\,M$ l'ens des monomorphismes, \deleted{\unsure{alors}}, et $\mathrm{Mon}_0$ le sous-ensemble de $\mathrm{Mon}$ formé des \deleted{\unsure{mono}} $i\in\mathrm{Mon}$ tels que source et but appartiennent à \deleted{\unsure{$M$}} $M_0$. On a alors \deleted{\unsure{ce}}

% slanted
\margin{\uline{NB} On a là un critère d'injectivité d'objets \deleted{\unsure{dans}} des catégories $M_{/Y}$ -- il suffit de tester la condition d'\unsure{inj.} par des $i$ mono. \unsure{tels que|i.e. que} $i$ soit dans $\mathrm{Mon}_0$.}

\[
\mathrm{Mon}\deleted{\unsure{=}}\subset\widetilde{\mathrm{Mon}_0},\quad \text{d'où }\left\{\begin{aligned}
\widetilde{\mathrm{Mon}}&=\widetilde{\mathrm{Mon}_0}\\
(\mathrm{Mon})_*&=(\mathrm{Mon}_0)_*\ \deleted{\unsure{(*)}}
\end{aligned}\right.\tag*{(7.1.2)}
\]

\deleted{et si il satisfait à la condition \textcircled{e} page 76 \unsure{lim.} \unsure{et de la lim} \unsure{filtr.} \unsure{de i mono est mono})}

\inserted{$\mathrm{Mon}$ est stable \deleted{\unsure{par}} cochangements de base, et \uline{$M_0$ est petite}}, et si \deleted{\illegible{1} \unsure{mêmes}}

\deleted{\[ \mathrm{Mon}=\widetilde{\mathrm{Mon}_0}\tag*{(7.1.3)}
\]}

% margin note written sideways
\margin{si $\mathrm{Mon}$ est stable par co-changement de base}

$2^\circ)$ Si \inserted{$M_0$ est \uline{petit} et si} les objets de $M_0$ sont accessibles \inserted{(p. ex. $M$ satisfait \textcircled{c} p.~76)}, et \deleted{si $M$ satisfait \textcircled{\unsure{e|c}} ($\mathrm{Mon}$ stable par \uline{$\lim$} filtrantes)}, alors il existe une factorisation fonctorielle des \deleted{tout} flèches $f$ de $M$ en

\[
f=pi,
\]

avec $i\in\mathrm{Mon}$, $p\in(\mathrm{Mon})_*$. (\uline{NB} Une flèche $p\colon X\longrightarrow Y$ de $M$ est dans $(\mathrm{Mon})_*$, sss elle fait de $X$ un objet \uline{injectif} de $M_{/Y}$ -- donc l'assertion de factorisation assure que dans les catégories induites $M_{/Y}$ ``il y a assez d'injectifs''). De plus, \deleted{\unsure{on}} on a

\[
\mathrm{Mon}=\widetilde{\mathrm{Mon}_0}\tag*{(7.1.3)}
\]

$3^\circ)$ Soit $W$ une partie de $\mathrm{Fl}(M)$ qui est \uline{$L$-accessible} \unsure{lire|i.e.} accessible, et stable par limites inductives filtrantes. [On suppose à présent que $M$ satisfait les conditions a) b) c) d) e) de p.~76, d'où une bonne théorie des parties $L$-accessibles de $\operatorname{Ob}M$ ou de $\mathrm{Fl}\,M$.] On suppose, \inserted{plus précisément,} \deleted{de plus} que $W$ satisfait les conditions
% archive sheet number: 212
\scanpage{214}

\mspage{88}

(7.1.4) $\left\{\vphantom{\dfrac{A}{A}}\right.$ a) $W$ est $L$-accessible (i.e. stable par $\varinjlim$ filtr., et accessible)

b) Soient $\alpha,\beta$ composables, avec $\unsure{\beta\alpha|\beta}$ et $\alpha\in W$, alors $\unsure{\beta|\alpha\beta|\beta\alpha}\in W$. (\uline{NB} Cette condition sera satisfaite si $W$ est un localiseur modérément saturé, et c'est à cette situation que nous allons appliquer notre énoncé).

c) $\mathrm{Mon}\cap W$ est stable par cochangement de base,

% margin note written sideways
\inserted{et que $M$ satisfait les conditions supplémentaires (7.1.5) $\left\{\vphantom{\dfrac{A}{A}}\right.$ e$'$) \unsure{les} \uline{$\lim$} filtrantes \unsure{commutent} au chgt. de base f) la condition d'exactitude (7.2.1) ci-dessous}

Soit $\Pi$ un cardinal tel que $W$ soit $\Pi$-accessible, et $M_0\subset\mathrm{Filt}^{\Pi}$. (donc $M_0$ essentiellement petite). Alors on a

\[
W\cap\mathrm{Mon}\subset\widetilde{(W\cap\mathrm{Mon})^\Pi}\tag*{(7.1.6)}
\]

(ici pour \deleted{tout} partie $C$ de $\mathrm{Fl}\,M$, on pose $C^\Pi=C\cap\mathrm{Fl}^{\Pi}M$), donc

\[
\widetilde{W\cap\mathrm{Mon}}=\widetilde{(W\cap\mathrm{Mon})^\Pi},\tag*{(7.1.7)}
\]

donc $\widetilde{W\cap\mathrm{Mon}}$ est $Q$-\uline{engendré} par une \uline{petite} partie de $W\cap\mathrm{Mon}$.

\uline{Corollaire} 7.2 : Soit $M$ satisfaisant a) b) c) d) e), et \unsure{aux} conditions supplémentaires \deleted{que si}

(7.2.1) $\left\{\vphantom{\dfrac{A}{A}}\right.$ \inserted{e$'$) les \uline{$\lim$} filtr. \deleted{\unsure{st}} commutent au chgt. de base} f) Si $A_1,A_2$ sont deux sous-objets d'un objet $B$ de $M$, \deleted{alors} $A_0=A_1\cap A_2$, alors $A_1\amalg_{A_0} A_2\hookrightarrow B$ est un monomorphisme

% slanted
\margin{Une autre façon de le dire : Si $C=\mathop{\mathrm{Sup}}\limits_{\mathrm{s.o.}(B)}(A_1,A_2)$, \deleted{\unsure{alors}} \unsure{ce carré cart.|carré cart.} \diagram{\begin{tikzcd}[ampersand replacement=\&,column sep=tiny,row sep=tiny]
A_0 \arrow[r,dash] \arrow[d] \& A_1 \arrow[d,dash]\\
A_2 \arrow[r,dash] \& C
\end{tikzcd}} est aussi cocart.}

\deleted{(7.2.2) g) $\mathrm{Mon}$ est stable par \unsure{cochgt. de base}.} \deleted{Soit} \unsure{$W$} une partie de $\mathrm{Fl}\,M$, satisfaisant \deleted{localiseur mod. saturé \unsure{dans}} \deleted{$M$, tel que} conditions suivantes

(7.2.3) $\left\{\vphantom{\dfrac{A}{A}}\right.$ a) $W$ est un localiseur modérément saturé \deleted{\unsure{$M$}} b) $W$ est $L$-accessible, i.e. \inserted{Il est} stable par \uline{$\lim$} filtrantes, et est accessible

c) $\mathrm{Mon}\cap W$ \deleted{\unsure{stable}}, \inserted{stable par cochgt de base, i.e. si on a un} diagramme \inserted{\uline{cocartésien}}

\diagram{\begin{tikzcd}[ampersand replacement=\&,row sep=small]
 A \arrow[r,hook,"i"] \arrow[d,"\alpha"'] \arrow[dr,phantom,"\mathrm{coc.}"] \& B \arrow[d,"\beta"]\\
 A' \arrow[r,"i'"'] \& B'
\end{tikzcd}}

avec $i$ mono\inserted{ et $\in W$}, alors \deleted{$\alpha\in W$ implique $\beta\in W$.} \inserted{$i'$ est mono et $\in W$.}

% slanted
\margin{Cf plus bas pour une quatrième condition d) sur $W$, qui assure que la situation donne un triple de Quillen clos}
% archive sheet number: 213
\scanpage{215}

\mspage{89}

Sous ces conditions, \unsure{\ldots|on a}

\[
\underbrace{\mathrm{Mon}}_{C}\cap W=\widetilde{\mathrm{Mon}\cap W}\qquad
\begin{aligned}&\text{i.e. }\mathrm{Mon}\cap W\text{ est}\\&\quad Q\text{-clos à gauche},\end{aligned}
\tag*{(7.2.\deleted{3}4)}
\]

et

\[
\begin{aligned}
&\text{Il existe une petite partie }\mathrm{TC}_0\subset\mathrm{Mon}\cap W\\
&\text{telle que }\mathrm{Mon}\cap W=\widetilde{\mathrm{TC}_0}.
\end{aligned}\tag*{(7.2.\deleted{\unsure{4}}5)}
\]

[De façon plus précise, si $\Pi\geq\unsure{\Pi_0|k}$ est un cardinal tel que $W$ soit $\Pi$-accessible, on peut prendre \deleted{pour} $\mathrm{TC}_0=(\mathrm{Mon}\cap W)^{\Pi}$ (ou une partie \uline{petite} de $(\mathrm{Mon}\cap W)^{\unsure{\Pi|P}}$ ayant les \unsure{$=^{\mathrm{s}}$} classes d'isomorphie).] \deleted{\unsure{pour t}} Donc (th. 1 du \S\,1) le couple $(\mathrm{TC}=\mathrm{Mon}\cap W,\ F=(\mathrm{TC})_*)$ est un couple de Quillen clos, et il y a $=$ une factorisation \uline{fonctorielle} de toute flèche $f$ en

\[
f=pi\ ;\quad i\in\unsure{\mathrm{TC}|\mathrm{TNC}},\ p\unsure{\in|\vdash}F.\tag*{(7.2.\unsure{4'|4}.)}
\]

\deleted{Mieux}, Si on pose

\[
C=\mathrm{Mon},\ \mathrm{TC}=\mathrm{Mon}\cap W,\ F=(\mathrm{TC})_*,\ \mathrm{TF}=C_*,\tag*{(7.2.5.)}
\]

le triple $(W,C,F)$ est un triple de Quillen clos, sss on \deleted{\unsure{suppose que l'on a}} \inserted{a les conditions \inserted{suppl.} suivantes}

$\left\{\vphantom{\dfrac{A}{A}}\right.$ (7.2.6) \textcircled{g} \deleted{\illegible{1}} (sur $M$ seulemt) $\mathrm{Mon}$ stable par cochg. de base.

\[
\underset{\substack{\Vert\,\mathrm{def}\\(\mathrm{Mon})_*}}{\mathrm{TF}}\subset W\deleted{\unsure{\mathrm{univ}}} \tag*{(7.2.\deleted{\unsure{6}}7)}
\]

i.e. si \inserted{d)} \unsure{toute} $f\colon X\mathrel{-}Y$ qui fait de $X$ un objet \uline{injectif} de $M_{/Y}$, est dans $W$ (donc aussi dans $W_{\mathrm{univ}}$).
% archive sheet number: 214
\scanpage{216}

\mspage{90}

\uline{Démonstration de} 7.1. On peut donner un énoncé un peu plus délicat, en supposant \inserted{d'emblée} donné, non une partie $M_0$ de $\operatorname{Ob}M$, mais une partie

\[
\mathrm{Mon}_0\subset\mathrm{Mon},\tag*{(7.1)}
\]

satisfaisant la condition suivante

(7.2) Pour \unsure{tte|tt} \deleted{flèche} mono $i\colon A\mathrel{\unsure{-|\to}} B$ dans $M$, il existe \inserted{$i'\colon A'\mathrel{-}B'$} $i'\in\mathrm{Mon}_0$, et une flèche $i\deleted{_0}'\longrightarrow i$ dans \uline{$\mathrm{Fl}$} $M$

\diagram{\begin{tikzcd}[ampersand replacement=\&,row sep=small]
 A' \arrow[r,dash,"i'"] \arrow[d] \& B' \arrow[d]\\
 A \arrow[r,"i"] \& B
\end{tikzcd}}\hfill,

telle que la flèche correspondante \deleted{\unsure{à}}

\[
C=A\amalg_{A'}B'\longrightarrow B
\]

soit un mono, et que $B'\mathrel{\unsure{-|\to}}B$ ne se factorise pas par $A$.

C'est cela qui était visé par la condition (7.1.1) sur $M_0$, qui implique (7.2) sur $\mathrm{Mon}_0$ défini en termes de $M_0$ comme dit au $1^\circ$.

Dans ce qui suit, on suppose donnée une partie $\mathrm{Mon}_0\subset\mathrm{Fl}\,M$ satisfaisant \unsure{(2)|(1)}. En pratique, elle sera souvent déduite d'une partie $M_0$ de $\operatorname{Ob}M$, génératrice pour les monomorphismes, comme \unsure{explicité|expliqué} dans $1^\circ$), mais on n'a pas besoin du fait que $i'\longrightarrow i$ soit un mono (i.e. 7.1.1' est plus fort que ce qu'il nous faut). P.~ex. si $B'\in M_0$ et $B'\longrightarrow B$ est donné, ne se factorisant pas par $A$, on pourra prendre $A'=B'\times_B A$ et $i'\colon A'\longrightarrow B'$ (pour assurer $i'\in\mathrm{Mon}_0$, il faudra \unsure{encore} supposer $M_0$ stable par passage aux sous-objets), et il y a des chances
% archive sheet number: 215
\scanpage{217}

\mspage{91}

que $C\mathrel{\unsure{-|\longrightarrow}}B$ sera alors bel et bien mono...

$1^\circ)$ Il faut prouver

\[
\mathrm{Mon}\subset\widetilde{\mathrm{Mon}_0}
\]

moyennant (7.2.) (et la stabilité de $M$ par \uline{$\lim$} filtrantes). Cela va résulter directement du

\uline{Lemme 7.3.} \unsure{Soient|Soit} $M$ une catégorie stable par limites ind. filtrantes, et $C_0,C$ $(\subset\mathrm{Fl}\,M)$ des parties de $\mathrm{Fl}\,M$. On suppose \textcircled{a} $C$ formée de monos, et \textcircled{b} que la condition suivante (généralisant (7.2) ci-dessus) soit vérifiée

(7.3.1) Pour tte flèche $i\unsure{\in|\colon} A\hookrightarrow B$ dans $C$, \inserted{telle que $i$ ne soit un iso,} il existe $i'\colon A'\mathrel{\unsure{-|\to}}B'$ dans $\widetilde C_0$, et une flèche $i'\longrightarrow i$ dans \uline{$\mathrm{Fl}$} $M$, telle que \inserted{dans} \deleted{la flèche \unsure{correspondante}} diagramme

\diagram{\begin{tikzcd}[ampersand replacement=\&,column sep=small]
 A' \arrow[r,"i'"] \arrow[d] \& B' \arrow[d] \arrow[dr,dashed] \& {}\\
 A \arrow[r,"{\deleted{\unsure{i}}\alpha}"] \arrow[rr,bend right=18,dashed,"i"'] \& C=\unsure{A'|A}\amalg_{A'}B' \arrow[r,"\beta"] \& B
\end{tikzcd}}\hfill,

la flèche $\beta$ soit dans $C$, et que $B'\longrightarrow B$ ne se factorise pas par $A$ (ou, ce qui revient au $=$, que le sous-objet de $B$ défini par $C$\inserted{ et $\beta$} ne soit pas égal à celui défini par $A$ et $i$).

On suppose de plus \textcircled{c} $C$ stable par $\varinjlim$ filtrantes, et
% margin note written sideways
\inserted{\textcircled{d} que pour tt $B\in\operatorname{Ob}M$, l'ens. des sous-objets de $B$ est petit. (Condition des plus anodines \deleted{\unsure{bien}} C'est une condition \uline{sur $M$}, \unsure{n'impliquant} pas $C$)} Alors on a

\[
C\subset\widetilde C_0,\ \text{donc }\widetilde C\subset\widetilde C_0,\tag*{(7.3.2)}
\]

donc, si $C_0\subset C$, on conclut
% archive sheet number: 216
\scanpage{218}

\mspage{92}

\[
\widetilde C=\widetilde C_0\tag*{(7.3.3)}
\]

(Ce lemme implique aussitôt la première \unsure{assertion} de 7.1. $1^\circ$.) Prouvons ce lemme, \deleted{\unsure{ou}} \inserted{i.e. que \deleted{tout} $i\in C$ est dans $\widetilde C_0$,} en construisant par récurrence transfinie \inserted{une suite} \deleted{des} \inserted{\unsure{str.|str} croissante de} sous-objets de $B$,

\[
A_\alpha\overset{j_\alpha}{\lhook\joinrel\relbar} B,\qquad \text{avec \inserted{les} }j_\alpha\in C,\ \text{et}
\]

\inserted{avec} \deleted{de telle façon que} \textcircled{$1^\circ$} $A_0=A$, \deleted{et} \inserted{On prouvera que} $A\xrightarrow{i_\alpha=i_{\alpha,0}}A_\alpha$ dans $\widetilde C_0$, \deleted{\unsure{et de}} et \deleted{Si $j$ est un iso} il y aura un ordinal final $\alpha$ tel que $A_\alpha=B$; \deleted{\unsure{et}} on en conclura donc que $i=i_\alpha$ est dans $\widetilde C_0$.

Si $i$ est un iso, on sait que $i\in\widetilde C_0$, il \deleted{\unsure{On doit}} n'y a rien à prouver (la constr. transfinie s'arrête au cran 0). Les \deleted{deux} \inserted{deux} pas de la construction, après le pas initial $1^\circ$, sont :

\textcircled{$2^\circ$} Soit $A_\alpha$ construit. Si $A_\alpha=B$, la construction est arrêtée. Si $A_\alpha\xrightarrow{j_\alpha} B$ \unsure{non iso|un iso}, on applique 7.3.1. à $j_\alpha\in C$, pour trouver un \deleted{\unsure{$-$}} $C\overset{j}{\hookrightarrow}B$, avec $j\in C$. et $C\supsetneqq A_\alpha$. De plus,

\[
i_{\alpha+1,\alpha}\colon A_\alpha\mathrel{\unsure{-|\to}} A_{\alpha+1}\overset{\mathrm{def}}=C\ \text{est dans }\widetilde C_0,\tag*{(7.3.4.)}
\]

car déduit de $i'\in\widetilde C_0$ par cochgt. de base.

\textcircled{$3^\circ$} Si $\alpha$ est ordinal limite, on pose

\[
A_\alpha=\varinjlim_{\beta<\alpha} A_\beta\tag*{(7.3.5)}
\]
% archive sheet number: 217
\scanpage{219}

\mspage{93}

et la flèche can.

\[
A_\alpha\xrightarrow{j_\alpha}B
\]

est dans $C$, car les $j_\beta$ le sont et $C$ est stable par limites \inserted{ind.} filtrantes. \deleted{On a \unsure{ainsi}}

On a ainsi achevé la description du système transfini. Pour \deleted{\unsure{tt}} $\alpha$, on a

\[
\mathrm{card}\,\alpha\leq\mathrm{card}\,\unsure{\mathrm{Ss}|\mathrm{Sos}}(B),
\]

puisque l'application $\beta\longmapsto A_\beta$ \deleted{est \unsure{de}} de l'ens $I_\alpha=\{\beta\mid\beta<\alpha\}$ dans $\mathrm{Ss}(B)$ est \uline{str} croissante. Donc il existe un ordinal $\alpha$ tel que $\mathrm{card}\,\alpha\leq\mathrm{card}\,\mathrm{Ss}(B)$, tel que $A_\alpha=B$. \deleted{Pour prouver que $j\in$} Mais alors $i\in\widetilde C_0$, comme limite inductive \inserted{(ou composée)} transfinie de flèches $\in\widetilde C_0$. (Il suffit de la condition $\forall\beta,\ i_{\beta+1,\beta}\in\widetilde C_0$ pour impliquer que tous les morphismes de transition $j_{\beta',\beta}$ sont dans $\widetilde C_0$.)

Revenons à 7.1, \deleted{\unsure{1°}} \inserted{$2^\circ)$}. Il reste à prouver que si $\mathrm{Mon}$ est stable par cochgts. de base, alors on a \deleted{$\mathrm{Mon}$}

\[
\mathrm{Mon}=\widetilde{\mathrm{Mon}_0}.
\]

Mais comme $\mathrm{Mon}$ est \inserted{déjà} stable par composition, et par limites inductives filtrantes \unsure{quelc.|finales|finies}, \deleted{\unsure{et cofacteurs}} \inserted{(donc aussi par facteurs directs)}, on voit par le cor. 2 du th. 1, \S\,1 \inserted{(p.~\unsure{19|18})} que $\mathrm{Mon}=\widetilde{\mathrm{Mon}_0}$. L'énoncé de factorisation est le th 1 du \S\,1.
% archive sheet number: 218
\scanpage{220}

\mspage{94}

Prouvons 7.1. $3^\circ)$, \deleted{\unsure{pour ce}} i.e. la formule

\[
W\cap\mathrm{Mon}\subset\widetilde{(W\cap\mathrm{Mon})^\Pi}.\tag*{(7.1.\deleted{\unsure{5}}6)}
\]

Pour ceci, nous \unsure{appliquons|appliquerons} encore le lemme 7.3. à

\[
C_0=(W\cap\mathrm{Mon})^\Pi,\qquad C=W\cap\mathrm{Mon}.
\]

Les conditions a), b), c) \inserted{d)} du lemme sont vérifiées. C'est clair pour toutes, sauf pour b), qui demande une vérification non tautologique. Soit donc

\[
f\colon A\hookrightarrow B,\qquad \begin{aligned}&f\in\mathrm{Mon}\cap\unsure{W'|W},\\&\inserted{f\text{ non iso.}}\end{aligned}
\]

On considère $B$ comme limite inductive

\[
B=\varinjlim_I B_i,\qquad I=\unsure{\mathrm{Ss}|\mathrm{Sos}}_\Pi(\deleted{\unsure{M}}B).
\]

% slanted
\margin{\uline{système inductif} \uline{$\Pi$-adapté}}

Pour tt $i$, on pose

\[
A_i=A\cap B_i=A\times_B B_i,
\]

et

\[
f_i\colon A_i\hookrightarrow B_i,
\]

donc

\[
f=\varinjlim_{i\in I} f_i.
\]

On a besoin de savoir que le syst. inductif des flèches $f_i$ est $\Pi$-adapté, ce qui nous amène à postuler la condition d'exactitude \inserted{suppl.} suivante sur $M$ :
% archive sheet number: 219
\scanpage{221}

\mspage{95}

(7.4.) e$'$) les \uline{$\lim$} filtrantes dans $M$ commutent au chgt de base, il suffit \unsure{=|même} de le supposer dans le cas très particulier \unsure{qu'on a} envisagé ici. \unsure{//}

Comme \deleted{$C$ \inserted{\deleted{$=\mathrm{Mon}$}} et $W$ sont $\Pi$-} $W$ est $\Pi$-accessible, et $(f_i)_{i\in I}$ est $\Pi$-adapté \deleted{\unsure{avec}}, comme $\lim f_i\in W$, il s'ensuit que l'ens \deleted{des} $I_W$ des $i\in I$ tels que $f_i\in W$ est cofinal. Comme $A\overset{f}{\hookrightarrow}B$ n'est pas iso, il existe donc un $i\in I_W$ tel que $B_i\hookrightarrow B$ ne se factorise pas par $A$. On a donc le diagramme

\diagram{\begin{tikzcd}[ampersand replacement=\&,column sep=small,row sep=small]
 A_i \arrow[r,hook,dash,"f_i"] \arrow[d,hook] \arrow[dr,phantom,"\mathrm{cocart}"] \& B_i \arrow[d,hook] \arrow[dr,dashed,"\mathrm{incl.}"] \& {}\\
 A \arrow[r,hook,"\alpha_i"] \arrow[rr,bend right=19,dashed,"f"'] \& C_i \arrow[r,hook,"\beta_i"] \& B
\end{tikzcd}}.

\deleted{On} avec $f_i\in(C\cap W)^\Pi\inserted{=C_0}$, donc $\alpha_i\in\widetilde C_0$; \deleted{\unsure{et}} On \deleted{\illegible{3}} \inserted{a besoin} de savoir que $\alpha_i\in W$, donc une condition supplémentaire sur $W$ :

(7.5.) Si on a un carré cocartésien

\diagram{\begin{tikzcd}[ampersand replacement=\&,column sep=small,row sep=small]
 A \arrow[r,hook,"i"] \arrow[d] \& B \arrow[d]\\
 A' \arrow[r,hook,"i'"] \& B'
\end{tikzcd}}\hfill,

avec $i\in\mathrm{Mon}\cap W$, alors $i'\in W$

% slanted
\margin{\uline{NB} Cette condition est vérifiée}

(j'ai renforcé (\unsure{7|4}.1.4.) en conséquence)\unsure{.}
% archive sheet number: 220
\scanpage{222}

\mspage{96}

Comme par hypothèse $f=\beta_i\alpha_i$ est dans $W$, et $\alpha_i\in W$, il s'ensuit par 7.1.4. a) que $\beta_i\in W$. Il faudra que nous sachions de plus que $\beta_i$ est mono, donc $\beta_i\in\mathrm{Mon}\cap W=C$, ce qui achèvera \deleted{de vérifier} d'établir la condition manquante b) de 7.3. C'est la condition \deleted{(7.1.1.)}\inserted{d'exactitude} qu'il me faut visiblement rajouter aux hypothèses \deleted{du} 7.1. 3°). Ayant fait cela, le lemme 7.3. s'applique et établit 7.1.3°.

\uline{Démonstration de} 7.2. Posons

\[
C=\mathrm{Mon},\quad\mathrm{TC}=C\cap W,\quad\mathrm{TC}_0=\mathrm{TC}^{\Pi},
\]

on a donc par 7.1. 3°, on a

\[
\underbrace{\mathrm{TC}_0}_{\substack{\text{ess}^{\mathrm{t}}\\\text{petite}}}\subset\mathrm{TC}\subset\widetilde{\mathrm{TC}_0}=\widetilde{\mathrm{TC}}
\]

et par le th. 1 \S1 (p. \deleted{\unsure{10}}12), cor. \unsure{1|2}, pour vérifier que $\mathrm{TC}=\widetilde{\mathrm{TC}_0}$ (i.e $\mathrm{TC}$ est $Q$-clos \unsure{à|:} g.) il \deleted{\unsure{sert}} \unsure{revient} au $=$ de voir qu'il est Q-stable, i.e \unsure{S|S°.}) être stable par \inserted{compositions, par} cochgt \deleted{de} \unsure{base|b.}, par \uline{$\lim$} ordinales strictes, par fact. directe. Or les deux dernières sont induites par le fait que $\inserted{\mathrm{TC}=}C\cap W$ est stable par \uline{$\lim$} filtrantes, vu que $W$ l'est, et $C$ aussi par la condition \textcircled{e} \inserted{(p. 76)} sur $M$. D'autre part, $C=\mathrm{Mon}$ et $W$ étant stables par composition, idem pour $C\cap W$. Donc il reste le cochgt. de
% archive sheet number: 221
\scanpage{223}

\mspage{97}

base, et c'est assuré par 7.2.3. c) (dont on avait déjà besoin dans 7.1., 3°, pour établir $C\cap W\subset\deleted{\unsure{C_{\mathrm{stab}}}}\widetilde{\mathrm{TC}_0}$. (\uline{NB} les hypothèses de 7.2. sont pratiquement celles de 7.1, 3°, sauf qu'on suppose que $W$ est un localiseur (\unsure{faiblement|au sens} : modérément saturé), de façon précise on n'a utilisé, jusqu'à présent, comme hyp. suppl., que la stabilité de $W$ par composition.)

\deleted{Si on rajoute les hypothèses \unsure{a) à|a) et}}

Il reste à établir que $(W,C=\mathrm{Mon.},F=(\mathrm{TC})_*)$ est un triple de Quillen clos, ssi les conditions 7.2.6 et (7.2.7), i.e.

\[
\left\{\begin{aligned}
&C\inserted{=\mathrm{Mon}.}\ \text{stable par cochangement de base}\\
&\mathrm{TF}=C_*=(\mathrm{Mon})_*\subset W
\end{aligned}\right.\tag*{(7.2.7)}
\]

sont satisfaites. La nécessité est évidente, par définition des triples de Quillen clos, et pour la suffisance, on note qu'il résulte du th. 7.1, \unsure{2°|3°} et de \deleted{la} partie déjà prouvée de 7.2, que $(C,\mathrm{TF}),(\mathrm{TC},F)$ sont des couples de Quillen clos, et on a de plus

\[
\mathrm{TC}=C\cap W,\quad \mathrm{TF}\subset F\cap W
\]

(puisque $\mathrm{TF}\subset W$ par 7.1.7, et $\mathrm{TF}\subset\unsure{F|R}$ \deleted{\unsure{puisque}} $\mathrm{TC}\subset C$). \deleted{Il s'ensuit} On sait
% archive sheet number: 222
\scanpage{224}

\mspage{98}

(prop. 4.3, cor. 4, page 36) qu'il en résulte que $(W,C,F)$ est un triple de Quillen clos. Cela achève la dém. de 7.1 et 7.2.

Résumons l'essentiel :

% slanted
\margin{\uline{NB} C'est \unsure{finalement|surtout} ce corollaire qui est le résultat principal de la présente section.}

\uline{Cor.} 7.4. Soit $M$ une catégorie satisfaisant les conditions a) b) c) d) e) (p. 76) et e$'$), f) ((7.2.1.) p. 88) enfin g) (7.2.6) i.e. stabilité de l'ens $\mathrm{Mon}$ des monomorphismes de $M$ par cochgt. de base. Soit $W$ \deleted{\unsure{dans}} $\subset\mathrm{Fl}\,M$ un localiseur dans $M$, qui est \uline{stable par} \uline{$\lim$} \uline{filtrantes} et de plus (condition plus subtile, mais \inserted{pratiquement} anodine) \uline{accessible}.

On pose

\[
C=\mathrm{Mono},\ \mathrm{TC}=C\cap W,\ \mathrm{TF}=C_*,\ F=(\mathrm{TC})_*.
\]

Alors : 1°) $(C,\mathrm{TF})$ est un couple de Quillen clos (énoncé indépendant de la donnée de $W$) et $=$ un couple de Quillen à factorisation fonctorielle.

% slanted
\margin{\uline{NB} Dans les cas 1° et 2°, $C,\mathrm{TC}=C\cap W$ sont Q-engendrés chacun par une \uline{petite} partie.}

2°) Pour \inserted{que $(\mathrm{TC},F)$ soit un couple} \deleted{$\mathrm{TC}$ est stable par cochangement de} de Quillen, il f. et s. que $\mathrm{TC}$ soit stable par cochangement de base. (Et alors il y a factorisation fonctorielle.)

3°) Pour que $(W,C,F)$ soit un triple de Quillen clos, il f. et s. que les conditions suivantes soient satisfaites

\[
\left\{\begin{aligned}
&\text{a) }\mathrm{TC}\overset{\mathrm{def}}{=}W\cap\mathrm{Mono}\quad\text{stable par cochgt. de base}\\
&\text{b) }\mathrm{TF}\overset{\mathrm{def}}{=}(\mathrm{Mono})_*=\text{ applications de \unsure{type} \unsure{injectif}}\\
&\qquad\subset W.
\end{aligned}\right.\tag*{(7.4.1)}
\]
% archive sheet number: 223
\scanpage{225}

\mspage{99}

4°) Pour que $(W,C,F)$ soit un triple de Quillen clos et propre à gauche (i.e. $C\subset\mathrm{Cof}_W$) il faut et il suffit qu'on ait la condition (7.4.1.) b) (i.e $W$ contient les morphismes de type injectif), et de plus

\[
C\subset\mathrm{Cof}_W,\quad \unsure{\text{ou}|\text{i.e.}}\inserted{\text{encore :}}\ \text{pour tout carré cocartésien}\tag*{(7.4.2)}
\]

\diagram{\begin{tikzcd}[ampersand replacement=\&,column sep=small,row sep=small]
A \arrow[r,"i"] \arrow[d,dash,"\alpha"'] \& B \arrow[d,dash,"\beta"]\\
A' \arrow[r,"i'"] \& B'
\end{tikzcd}},

avec $i$ un mono, alors $\alpha\in W\Longrightarrow\beta\in W$.

5°) Pour que $(W,C,F)$ soit un triple de Quillen clos et \uline{propre}, il f et s. que les deux conditions de 4°) soient satisfaites, ainsi que l'on ait

\[
F\subset\mathrm{Fib}_W\quad\unsure{\text{ou}|\text{\deleted{ssi}}}\ \text{encore : pour tout carré cartésien}\tag*{(7.4.3)}
\]

\diagram{\begin{tikzcd}[ampersand replacement=\&,column sep=small,row sep=small]
X \arrow[d,dash,"f"'] \& X' \arrow[l,"q"'] \arrow[d,dash,"f'"]\\
Y \& Y' \arrow[l,dash,"p"]
\end{tikzcd}}, si

$f\in F$ (i.e. $f\in(C\cap W)_*$) et $p\in W$, on a $q\in W$.

Dém de 7.4. Les parties 1°, 2°) 3°) ne sont qu'un résumé de ce qui était établi. Pour 4°, il faut voir que les conditions dites impliquent (7.4.1.) a), Or on \deleted{qui est une conséquence formelle du fait qu'on sait déjà que $(C,\mathrm{TF})$ est un couple de Quillen, et \unsure{i.e.}} $C\subset\mathrm{Cof}_W$, donc $C\cap W\subset\mathrm{Cof}_W\cap W\deleted{\unsure{=|{\subset}}\unsure{\mathrm{Cof}}}=\unsure{W}^{\mathrm{univ}}$ \deleted{\unsure{et}}
% archive sheet number: 224
\scanpage{226}

\mspage{100.}

donc $C\cap W=C\cap W^{\mathrm{univ}}$, \deleted{\unsure{qui}} \unsure{qui} est stable par cochgt. de base puisque $C$ et $W^{\mathrm{univ}}$ le sont. Quant à 5°) c'est une tautologie, \deleted{compte} tenu de 4°.

8. \uline{Application aux catégories de modèles} $(\widehat A,W_A)$.

Soit $A$ dans $\mathrm{Cat}$, et considérons dans $\widehat A$ le localiseur $W_A$, image inverse par

\[
\left\{\begin{aligned}i_A:\quad&\widehat A\longrightarrow\mathrm{Cat}\\&F\longmapsto A/F\end{aligned}\right.\tag*{(8.1.)}
\]

d'un localiseur fondamental donné $W$, satisfaisant les conditions

\[
W\bigl(\underset{\text{sat.}}{1},\underset{\text{loc.}}{2}, \underset{\substack{\text{\unsure{dépl.}}\\\text{fin.}}}{3}, \underset{\substack{\text{\unsure{conn}}\\\Pi_0}}{5}, \underset{\varinjlim}{6}, \underset{\substack{\text{carrés}\\\text{\unsure{coc.} en}\\\text{``M.V.''}}}{7\,\mathrm{bis}\deleted{\unsure{3}}}, \underset{\text{acc.}}{8}.\bigr)\tag*{(8.2.)}
\]

\margin{où $W(8)$ est l'accessibilité de $W$ (donc $W$ est L-accessible s'il satisfait aussi $W(\unsure{6|5})$).}

\deleted{\unsure{Comme} $\widehat A$ \unsure{est}} Pour pouvoir \unsure{à présent|maintenant} appliquer les développements sur les parties accessibles d'une grosse catégorie, \deleted{il faut} du \S\,5, il faut vérifier les conditions a) b) c) d) e) \deleted{\unsure{\S}} du \S\,5, tant \deleted{dans} $\mathrm{Cat}$, que (surtout) dans $\widehat A$.

a) Stabilité par \deleted{\unsure{petites}} \uline{$\lim$} filtrantes, OK.

b) Filtration cardinale. Bien connue pour $\widehat A=\underline{\mathrm{Hom}}\,(A,\mathrm{Ens})$, à partir de celle de $\mathrm{Ens}$ par les cardinaux [cf SGA 4 I 9.2\unsure{4|.4}, p. 172]. Pour $\mathrm{Cat}$, on définit une filtration cardinale

\[
(\mathrm{Cat}^{\Pi})_{\Pi\geq\Pi_0},\quad\text{où }\Pi_0=\text{Cardinal des ens. dénombrables}\tag*{(8.3)}
\]
% archive sheet number: 225
\scanpage{227}

\mspage{101}

Par la condition, pour $A$ dans $\mathrm{Cat}$

\[
A\in\mathrm{Filt}^{\Pi}(\mathrm{Cat})=\mathrm{Cat}^{\Pi}\quad\Longleftrightarrow\quad\mathrm{card}\,\mathrm{Fl}(A)\leq\Pi\tag*{(8.4)}
\]

Il faut vérifier \deleted{\unsure{des}} les conditions a) b) c) de loc. cit 9.12. La condition a) est évidente, b) est contenue dans la condition plus forte de stabilité de $\mathrm{Cat}^{\Pi}$ par $\varinjlim_I$, pour $I$ une cat. telle que $\mathrm{Card}\,\mathrm{Fl}\,I\leq\Pi$ (c'est la condition d$_{\unsure{\alpha|0}}$) de ces notes, p. 45-46). Notons que $\mathrm{Cat}^{\Pi}$ est également stable par $\varprojlim$ finies, (et si $\Pi$ est de la forme $2^c$, il est stable $=$ par limites projectives $\varprojlim_I$, avec $\mathrm{card}\,\mathrm{Fl}\,I\leq c$ \ldots). La condition c) se vérifie aisément par la méthode du \S\,5, en notant que pour $X$ dans $\mathrm{Cat}$, l'ens

\[
\unsure{\mathrm{Ss}|\mathrm{Sos}}_{\Pi}(X)
\]

des sous-catégories de $X$ qui sont dans $\mathrm{Cat}^{\Pi}$ est filtrant croissant et grand devant $\Pi$, et que $X$ en est la limite inductive. D'autre part, il est clair aussi que si $X\in\mathrm{Cat}^{\Pi'}$ avec $\Pi'>\Pi$, alors \deleted{\unsure{Card}}

\[
\mathrm{Card}\,\mathrm{Ss}_{\Pi}(X)\leq{\Pi'}^{\Pi}
\]

ce qui achève d'établir la condition c).

De plus, la condition d) de p. 45 se vérifie immédiatement \deleted{pour} les $\mathrm{Cat}^{\Pi}$, et aussi pour les $\widehat A^{\Pi}$ (pour $\Pi\geq\mathrm{Card}\,\mathrm{Fl}\,A$).

Ainsi que les mono sont stables par $\varinjlim$ filtrantes (condition e) page 46).
% archive sheet number: 226
